\section{Results}
\label{sec:results}

We evaluate token pruning across two VLA architectures on the LIBERO-Object benchmark~\cite{liu2023libero}, comprising 10 tasks with 50 episodes each (500 total). All success rates are reported with Wilson score 95\% confidence intervals.

\subsection{SmolVLA: Cliff-Edge Under All Pruning Methods}
\label{sec:smolvla_results}

SmolVLA~\cite{shukor2025smolvla} uses a pixel-shuffle vision connector that spatially rearranges patch embeddings before projecting them into the language model's embedding space. This operation creates inter-token spatial dependencies: each output token encodes information from a local grid of input patches. We hypothesize that this coupling makes the token sequence fragile under any form of pruning, since removing or masking a single token disrupts the spatial structure that downstream layers expect.

\paragraph{Setup.}
We use a SmolVLA checkpoint fine-tuned on LIBERO-Object with the ACIS (Adaptive Contextual Importance Scoring) module attached but pruning disabled during training. The baseline (bypass mode) achieves 44.8\% SR (224/500, CI [40.5\%, 49.2\%]). We then activate pruning under six configurations that vary three axes: scoring function (learned ACIS vs.\ random), retention ratio $k \in \{0.5, 0.7, 0.9\}$, and pruning mechanism (zero-masking, true token removal, token merging via ToMe~\cite{bolya2023tome}).

\paragraph{All configurations collapse to 0\% SR.}
Table~\ref{tab:smolvla_ablation} reports the full ablation. Every pruning configuration produces 0\% success rate across all 500 episodes (Wilson CI [0\%, 0.76\%]). Three observations make this result striking:

\begin{enumerate}
    \item \textbf{Scoring-independent.} Random token selection performs identically to the learned ACIS scorer. The failure does not stem from poor importance estimation.
    \item \textbf{Ratio-independent.} Retaining 90\% of tokens ($k{=}0.9$, removing only 10\%) causes the same catastrophic failure as retaining 50\%. There is no graceful degradation at any operating point.
    \item \textbf{Mechanism-independent.} Zero-masking, true removal (sequence shortening), and token merging all produce identical 0\% SR. The failure is not an artifact of a particular implementation choice.
\end{enumerate}

The convergence of all six conditions to 0\% SR, against a 44.8\% baseline, constitutes a cliff-edge failure pattern: performance does not degrade gradually but collapses entirely once any tokens are disturbed.

\begin{table}[t]
\centering
\small
\caption{SmolVLA token pruning ablation on LIBERO-Object. All six pruning configurations collapse to 0\% SR regardless of scoring function, retention ratio, or mechanism. The baseline achieves 44.8\% with pruning disabled.}
\label{tab:smolvla_ablation}
\begin{tabular}{llccrr}
\toprule
Scoring & $k$ & Mechanism & SR (\%) & Episodes & Wilson 95\% CI \\
\midrule
\multicolumn{2}{l}{\textit{Baseline (bypass)}} & -- & 44.8 & 500 & [40.5, 49.2] \\
\midrule
ACIS & 0.5 & Zero-mask & 0.0 & 500 & [0, 0.76] \\
ACIS & 0.7 & Zero-mask & 0.0 & 500 & [0, 0.76] \\
ACIS & 0.9 & Zero-mask & 0.0 & 500 & [0, 0.76] \\
Random & 0.5 & Zero-mask & 0.0 & 500 & [0, 0.76] \\
ACIS & 0.5 & True removal & 0.0 & 500 & [0, 0.80] \\
ACIS & 0.5 & Merge (ToMe) & 0.0 & 500 & [0, 0.76] \\
\bottomrule
\end{tabular}
\end{table}

\paragraph{Geometric prior confirms the pattern.}
As additional evidence, we evaluate a non-learned geometric prior (Sobel edge-based scoring) at $k \in \{0.5, 0.7, 0.9\}$, all producing 0\% SR.\footnote{The geometric prior experiments use a different model checkpoint than the ACIS bypass baseline, so we report them separately from Table~\ref{tab:smolvla_ablation}. The qualitative conclusion is identical: any form of token disturbance causes complete failure.}

\paragraph{Cross-suite confirmation.}
The cliff-edge pattern generalizes beyond LIBERO-Object. On LIBERO-Spatial (baseline 51.8\% SR) and LIBERO-Goal (baseline 63.0\% SR), ACIS pruning at $k{=}0.5$ with true removal produces 0\% SR in both suites (0/500 episodes each, CI [0\%, 0.76\%]). The failure is not task-specific but architectural.

\paragraph{Interpretation.}
The pixel-shuffle connector in SmolVLA creates a structured token lattice where spatial adjacency carries semantic meaning. When any subset of tokens is removed, masked, or merged, the resulting sequence violates the implicit spatial contract that downstream transformer layers learned during pretraining. This interpretation aligns with the connector's design: pixel-shuffle compresses spatial resolution by interleaving channels, producing tokens whose meaning depends on their position within the original grid. We formalize this analysis in Section~\ref{sec:analysis}.


\subsection{OpenVLA: Graceful Degradation}
\label{sec:openvla_results}

OpenVLA~\cite{kim2024openvla} uses a two-layer MLP projector to map vision encoder outputs (from DINOv2 and SigLIP~\cite{zhai2023siglip}) into the Llama-2 language model's embedding space. Unlike pixel-shuffle, the MLP projector processes each patch token independently through a shared linear transformation. This design preserves per-token independence: the representation of one token carries no implicit dependency on neighboring tokens' presence.

\paragraph{Setup.}
We evaluate the OpenVLA-7B checkpoint on LIBERO-Object under random token removal (true removal, no learned scorer). The baseline achieves 74.0\% SR (370/500 episodes). We apply pruning at $k{=}0.5$, retaining 50\% of vision tokens selected uniformly at random.

\paragraph{Graded degradation instead of collapse.}
Table~\ref{tab:openvla_degradation} shows the results. At $k{=}0.5$, OpenVLA achieves 37.4\% SR (187/500), a relative drop of 49.5\% from baseline. This stands in sharp contrast to SmolVLA, where the same 50\% retention ratio produces 0\% SR. OpenVLA loses performance proportionally to the information removed, rather than collapsing catastrophically.

\begin{table}[t]
\centering
\small
\caption{OpenVLA token pruning on LIBERO-Object. Unlike SmolVLA, OpenVLA shows graded degradation under random token removal, consistent with its MLP projector preserving per-token independence.}
\label{tab:openvla_degradation}
\begin{tabular}{cccr}
\toprule
$k$ & SR (\%) & Relative Drop (\%) & Episodes \\
\midrule
1.0 (baseline) & 74.0 & -- & 500 \\
0.5 & 37.4 & 49.5 & 500 \\
\bottomrule
\end{tabular}
\end{table}

\paragraph{Per-task analysis reveals graded variance.}
The per-task breakdown at $k{=}0.5$ shows a wide range of success rates: from 2\% (butter) to 72\% (orange\_juice). Tasks involving visually distinct target objects (orange\_juice, alphabet\_soup at 60\%) retain higher success rates, while tasks requiring fine-grained visual discrimination among similar objects (butter at 2\%, cream\_cheese at 8\%) suffer the most. This graded, task-dependent degradation pattern confirms that OpenVLA retains partial visual information under pruning rather than losing all task-relevant signal.

\paragraph{Architectural contrast.}
The divergence between SmolVLA and OpenVLA under identical pruning conditions (random selection, $k{=}0.5$, true removal) isolates the vision connector as the determining factor. Both models receive the same visual input and operate on the same benchmark. The only structural difference in the vision-to-language pathway is the connector: pixel-shuffle (SmolVLA) versus MLP projector (OpenVLA). SmolVLA drops from 44.8\% to 0\%; OpenVLA drops from 74.0\% to 37.4\%. This architecture-dependent compressibility gap motivates our analysis of token spatial structure in Section~\ref{sec:analysis}.
